\documentclass[10pt]{article}
\usepackage[margin=0.7in]{geometry}
\usepackage{amsmath,amssymb,amsthm,microtype}
\newtheorem{theorem}{Theorem}
\newcommand{\R}{\mathcal R}
\newcommand{\A}{\mathcal A}
\begin{document}
\begin{center}
{\Large Direct Merging and Contracting Returns for Affine Transfer}\par\medskip
Collatz proof project\quad---\quad September 5, 2026
\end{center}
Let $T$ be the shortcut Collatz map, $j_t(n)$ its odd-step count, and
$\R(n)$ the assertion that $n$ reaches one. Put
$\A(u):=[\R(3u+2)\Rightarrow\R(27u+20)]$.
The previous formalization proves that $\forall u\geq0,\A(u)$ is
equivalent to full Collatz. Here we give direct and inductive certificates
on particular arithmetic progressions; the universal statement remains open.

\begin{theorem}[Direct merging]
For every natural $Q$, $\A(36+64Q)$ holds.
\end{theorem}
\begin{proof}
At the base parameter, $T^6(110)=T^6(992)=47$, with odd counts three
and one. Exact binary-cylinder transport therefore gives
\[
T^6(110+192Q)=T^6(992+1728Q)=47+81Q.
\]
Reaching one transfers through the common endpoint in either direction.
Lean checks both the base computations and the extension to every $Q$.
\end{proof}

\begin{theorem}[Contracting parameter return]
For every natural $Q$, set $u=2308+4096Q$ and $v=45+81Q$. Then
\[
v<u,\qquad T^{12}(3u+2)=3v+2,\qquad
T^{12}(27u+20)=27v+20.
\]
Consequently $\A(v)\Rightarrow\A(u)$. The premise $\A(v)$ is retained;
the theorem does not prove transfer unconditionally on this progression.
\end{theorem}
\begin{proof}
The base endpoints are $T^{12}(6926)=137$ and $T^{12}(62336)=1235$,
both with four odd steps. Parameter increments of $4096Q$ give endpoint
increments $243Q$ and $2187Q$, respectively. These are precisely the
increments obtained by replacing the parameter $45$ with $45+81Q$.
Also $u-v=2263+4015Q>0$.
If $3u+2$ reaches one, so does its shifted value $3v+2$. Apply $\A(v)$,
then transport convergence back from $27v+20$ to $27u+20$.
Every step is kernel checked in \texttt{Collatz.AffinePairReturn}.
\end{proof}

\paragraph{Exact search criterion.}
On the cylinder $u=r+2^tQ$, the two endpoints are
\[
x+3^{j+1}Q,\qquad y+3^{k+3}Q,
\]
where $x=T^t(3r+2)$, $y=T^t(27r+20)$ and $j,k$ are the respective
odd counts. Direct merging requires $x=y$ and $j=k+2$.
A return requires $y=9x+2$, $j=k$ and $x\equiv2\pmod3$.
Then $v=(x-2)/3+3^jQ$; checking $(x-2)/3<r$ and $3^j\leq2^t$
proves $v<u$ for every quotient. Base endpoint coincidence alone is
insufficient for either certificate.

\paragraph{Finite evidence.}
The reproducible search stops extending a cylinder once either certificate
is found. Through depth 18, its 3,236 first certificates classify 30,176
of the $2^{18}$ parameter residues by direct merging and 1,267 by a
contracting return; 230,701 remain unresolved. The return count is
conditional induction coverage, not a count of independently proved
transfer statements. These census figures are Python evidence, not a
kernel completeness theorem. Tests compare the full small-depth census
against direct orbits and replay every saved certificate at $Q=0,1,7$.
The formulas above, rather than those sample quotients, justify lifting.

\paragraph{A boundary for aligned merging.}
Lean also proves that both 2 and 20 reach one but $T^t(2)\ne T^t(20)$
for every $t$. Their joint orbit enters the two-state cycle $(1,2),(2,1)$.
Thus parameter zero needs a base case or another argument, even though
its transfer is true. More generally, an identity on an infinite parameter
progression using fixed path lengths $s,t$ forces equality of endpoint
slopes. For these partners that equality is
$3^{j+1}/2^s=3^{k+3}/2^t$, hence $s=t$ and $j=k+2$.
The slope deduction is a written arithmetic observation, not an additional
Lean theorem. The single exceptional parameter does not refute a proof
combining base cases and contracting returns.

\paragraph{Remaining obligation and verification.}
A complete induction needs every parameter to be a base case, admit a
direct certificate, or return to a smaller parameter. No such coverage
theorem is proved. Six public declarations in \texttt{AffinePairReturn}
enter the axiom audit. The search and saved data are
\texttt{analysis/affine\_pair\_search.py} and its matching JSON file.
This work supplies induction steps, not the full induction, and makes
no mathematical priority claim.
\end{document}
